\documentclass{article}

% NeurIPS 2025 submission style.  The detailed proofs are preserved in the appendix.
\usepackage{neurips_2026}

\usepackage{amsmath,amssymb,amsthm}
\usepackage{algorithm}
\usepackage{algorithmic}
\usepackage{mathtools}
\usepackage{xcolor}
\usepackage{thmtools}
\usepackage{thm-restate}
\usepackage{booktabs}
\usepackage{array}
\usepackage{enumitem}
\usepackage{microtype}
\usepackage{hyperref}
\usepackage{url}


\setlist{nosep,leftmargin=*}
\setlength{\textfloatsep}{8pt plus 2pt minus 2pt}
\setlength{\floatsep}{7pt plus 2pt minus 2pt}
\setlength{\intextsep}{7pt plus 2pt minus 2pt}
\setlength{\abovedisplayskip}{6pt plus 2pt minus 2pt}
\setlength{\belowdisplayskip}{6pt plus 2pt minus 2pt}
\setlength{\abovedisplayshortskip}{4pt plus 2pt minus 2pt}
\setlength{\belowdisplayshortskip}{4pt plus 2pt minus 2pt}

% Theorem environments
\newtheorem{theorem}{Theorem}
\newtheorem{lemma}[theorem]{Lemma}
\newtheorem{proposition}[theorem]{Proposition}
\newtheorem{corollary}[theorem]{Corollary}
\newtheorem{definition}[theorem]{Definition}
\newtheorem{remark}[theorem]{Remark}
\newtheorem{assumption}[theorem]{Assumption}

% Shorthand
\newcommand{\R}{\mathbb{R}}
\newcommand{\E}{\mathbb{E}}
\newcommand{\Prob}{\mathbb{P}}
\newcommand{\Reff}{R_{\mathrm{eff}}}
\newcommand{\Vhat}{\widehat{V}}
\newcommand{\Qhat}{\widehat{Q}}
\newcommand{\Vtilde}{\widetilde{V}}
\newcommand{\muhat}{\widehat{\mu}}
\newcommand{\xrmark}[2]{\xrightarrow[#2]{#1}}

\title{Effective Resistance Exploration for Monte Carlo Tree Search}

\author{
  Tuan Dam \\
  SOICT\\
  HUST\\
  Hanoi, vietnam \\
  \texttt{tuandq@soict@hust.edu.vn}
}

\begin{document}
\maketitle

\begin{abstract}
We introduce MCTS-ER, a family of Monte-Carlo Tree Search algorithms augmented with an effective-resistance (ER) bonus computed from search-tree visit counts.  The ER term acts as a path-level uncertainty signal: an under-visited edge contributes a large resistance and makes the whole trajectory more exploratory.  We analyze ER in two settings.  First, for Stochastic-Power-UCT, adding $C_2/T_{s,a}$ to the polynomial exploration bonus preserves the $O(n^{-1/2})$ root-value convergence rate.  Second, for convex-regularized MCTS (MENTS, RENTS, TENTS), injecting $c_2/T_{s,a}$ into Q-values before the regularized policy is computed preserves the exponential convergence envelope $O(t\exp(-t/(\log t)^3))$.  In TENTS-ER, the ER perturbation also modulates the Tsallis active set, producing broad early exploration and sparse late exploitation.  Full proof bodies are refereed in Appendix~\ref{app:full}.
\end{abstract}

\section{Introduction}
\label{sec:main_intro}
Monte-Carlo Tree Search (MCTS) combines forward simulation, tree expansion, and value backup, and it is often driven by local exploration bonuses such as UCT.  Local bonuses, however, do not distinguish between a node reached through a well-supported trajectory and a node reached through a single weakly visited bottleneck edge.  We add a structural signal based on effective resistance: on a tree, the resistance from the root to a node is exactly the sum of edge resistances on the unique path.  With edge resistance $r(e)=1/N(e)$,
\begin{equation}
\Reff(s_0,s_h)=\sum_{e\in\mathrm{path}(s_0,s_h)}\frac{1}{N(e)}.
\label{eq:main_er_path}
\end{equation}
Thus a small visit count on any edge increases the uncertainty assigned to the whole trajectory.  At a fixed node $s_h$, the part of \eqref{eq:main_er_path} that differs across actions is simply $1/T_{s_h,a}$, so the global path signal induces a local additive selection bonus while retaining a tree-wise interpretation.

\paragraph{Contributions.}
We reorganize the results around two ER-augmented MCTS families.  For \emph{Power-UCT-ER}, we prove that the ER bonus is asymptotically dominated by the polynomial bonus and hence preserves the Stochastic-Power-UCT concentration recursion.  For \emph{E3W-ER} with maximum entropy, relative entropy, or Tsallis entropy regularization, we prove that the ER perturbation decays quickly enough to be absorbed by the existing exponential convergence analysis.  We additionally derive regret and value-error bounds for all three regularizers and identify the adaptive sparsity effect unique to TENTS-ER.  The main paper states the algorithms and results; Appendix~\ref{app:full} contains the original technical development and unchanged proofs.

\section{Preliminaries and Algorithms}
\label{sec:main_prelim}
We consider a discounted MDP $\mathcal{M}=\langle\mathcal{S},\mathcal{A},R,P,\gamma\rangle$ with finite action set and fixed planning depth $H$.  The depth-$H$ truncated optimal values are defined by
\begin{align}
\widetilde{Q}(s_h,a)&=r(s_h,a)+\gamma\sum_{s_{h+1}}P(s_{h+1}|s_h,a)\widetilde{V}(s_{h+1}),
& \widetilde{V}(s_h)&=\max_a\widetilde{Q}(s_h,a).
\label{eq:main_truncated_values}
\end{align}
For the power-mean backup, after $t$ trajectories the empirical estimates are
\begin{align}
\Vhat_t(s_h)&=\left(\sum_{a\in\mathcal{A}_{s_h}}\frac{T_{s_h,a}(t)}{t}\bigl(\Qhat_{T_{s_h,a}(t)}(s_h,a)\bigr)^p\right)^{1/p},
\label{eq:main_power_backup}\\
\Qhat_t(s_h,a)&=\frac{1}{t}\sum_{i=1}^{t}\left[r^i(s_h,a)+\gamma\Vhat_{T^{s_{h+1}}_{s_h,a}(i)}(s_{h+1})\right].
\label{eq:main_q_backup}
\end{align}
The analysis uses polynomial concentration: $\Vhat_n\xrightarrow{\alpha,\beta}V$ if there is $c>0$ such that, for all $\varepsilon>n^{-\alpha/\beta}$,
\begin{equation}
\Prob(|\Vhat_n-V|>\varepsilon)\leq c n^{-\alpha}\varepsilon^{-\beta}.
\label{eq:main_poly_conc}
\end{equation}

\begin{definition}[Tree effective resistance]
For a search tree $\mathcal{T}$ rooted at $s_0$, with edge visit counts $N(e)$, define
\[
\Reff(s_0,s_h)=\sum_{e\in\mathrm{path}(s_0,s_h)}\frac{1}{N(e)}.
\]
At node $s_h$, selecting action $a$ toward child $s_{h+1}$ changes this quantity by $1/T_{s_h,a}(t)$, because the prefix resistance $\Reff(s_0,s_h)$ is action-independent.
\end{definition}

\subsection{ER-Augmented Search Rules}
\label{sec:main_algorithms}
\paragraph{Power-UCT-ER.}
At node $s_h$, the selected action is
\begin{equation}
 a_h=\arg\max_{a\in\mathcal{A}_{s_h}}\left\{\Qhat_{T_{s_h,a}(t)}(s_h,a)+C_1\frac{T_{s_h}(t)^{b_{h+1}/\beta_{h+1}}}{T_{s_h,a}(t)^{\alpha_{h+1}/\beta_{h+1}}}+\frac{C_2}{T_{s_h,a}(t)}\right\}.
\label{eq:main_power_er_select}
\end{equation}
The first exploration term is the polynomial Stochastic-Power-UCT bonus; the second is the incremental ER contribution of edge $(s_h,a)$.  The value backup remains exactly the power mean in \eqref{eq:main_power_backup}.

\begin{algorithm}[t]
\caption{MCTS-ER main selection rule}
\label{alg:mcts-er-main}
\begin{algorithmic}[1]
\REQUIRE Root $s_0$, depth $H$, constants $C_1,C_2$, power $p$, rollout policy $\pi_0$
\FOR{$t=1,2,\ldots,n$}
    \STATE Traverse the tree; at depth $h$, choose $a_h$ by \eqref{eq:main_power_er_select}.
    \STATE Expand a leaf if depth is below $H$; evaluate it using $\pi_0$.
    \STATE Back up Q-values by \eqref{eq:main_q_backup} and values by \eqref{eq:main_power_backup}.
\ENDFOR
\RETURN $\hat a=\arg\max_a\Qhat_{T_{s_0,a}(n)}(s_0,a)$.
\end{algorithmic}
\end{algorithm}

\paragraph{Convex-regularized E3W-ER.}
Let $\Omega$ be a strongly convex regularizer and let $\Omega^*$ be its Fenchel dual.  The regularized value and policy are
\begin{equation}
 V_\Omega(s)=\Omega^*(Q_\Omega(s,\cdot)/\tau),\qquad
 \pi_\Omega(\cdot|s)=\nabla\Omega^*(Q_\Omega(s,\cdot)/\tau).
\label{eq:main_reg_policy}
\end{equation}
E3W-ER augments Q-values before evaluating the regularized policy:
\begin{equation}
 Q_t^{\mathrm{ER}}(s,a)=Q_t(s,a)+\frac{c_2}{T_{s,a}(t)},\qquad
 \pi_t^{\mathrm{ER}}(\cdot|s)=(1-\lambda_t)\nabla\Omega^*(Q_t^{\mathrm{ER}}(s,\cdot)/\tau)+\lambda_t\mathrm{Unif}(\mathcal{A}).
\label{eq:main_e3w_er}
\end{equation}
For maximum entropy, this gives a softmax over $Q+c_2/T$.  For relative entropy, the same perturbation appears inside the exponent relative to the previous policy.  For Tsallis entropy, the policy is sparse:
\begin{equation}
\pi_{\mathrm{TENTS\text{-}ER}}(a|s)=\left[\frac{Q(s,a)+c_2/T_{s,a}}{\tau}-\psi\left(\frac{Q(s,\cdot)+c_2/T_{s,\cdot}}{\tau}\right)\right]_+.
\label{eq:main_tents_er}
\end{equation}
Small $T_{s,a}$ can temporarily add an action to the Tsallis active set; as visits grow, this extra term vanishes and the policy contracts to the unaugmented sparse set.

\begin{table}[t]
\centering
\caption{ER insertion point and asymptotic effect.}
\label{tab:main_alg_summary}
\small
\begin{tabular}{p{0.22\linewidth}p{0.36\linewidth}p{0.32\linewidth}}
\toprule
Family & ER-augmented decision quantity & Preserved rate \\
\midrule
Power-UCT-ER & $\Qhat+C_1T^{b/\beta}/T_a^{\alpha/\beta}+C_2/T_a$ & $O(n^{-1/2})$ expected root payoff \\
MENTS-ER & $\mathrm{softmax}((Q+c_2/T)/\tau)$ & $\exp(-t/(\log t)^3)$ \\
RENTS-ER & relative-entropy policy on $(Q+c_2/T)/\tau$ & $\exp(-t/(\log t)^3)$ \\
TENTS-ER & sparse Tsallis map on $(Q+c_2/T)/\tau$ & $\exp(-t/(\log t)^3)$ with adaptive active set \\
\bottomrule
\end{tabular}
\end{table}

\section{Main Theoretical Results}
\label{sec:main_results}
\subsection{Power Mean MCTS-ER}
\label{sec:main_power_results}
The core technical observation is that ER decays faster than the polynomial bonus under the Stochastic-Power-UCT parameter conditions.  For parent visit count $n$ and child visit count $s$,
\begin{equation}
B(n,s)=C_1\frac{n^{b/\beta}}{s^{\alpha/\beta}}+\frac{C_2}{s}.
\label{eq:main_combined_bonus}
\end{equation}
When $\alpha/\beta<1$, the ratio $(C_2/s)/(C_1n^{b/\beta}s^{-\alpha/\beta})$ is proportional to $s^{\alpha/\beta-1}$ and tends to zero as $s\to\infty$.  The ER-modified suboptimal-arm threshold is therefore only a constant-shift perturbation of the original threshold:
\begin{equation}
A_a^{\mathrm{ER}}(n)\leq A_a(n)+\frac{4C_2}{\Delta_a},\qquad
A_a(n)=\left(\frac{2C_1}{\Delta_a}\right)^{\beta/\alpha}n^{b/\alpha}.
\label{eq:main_threshold}
\end{equation}
This lets every concentration lemma in the power-mean proof proceed with the same exponents and modified constants.

\begin{theorem}[Bandit concentration with ER]
\label{thm:main_bandit_er_statement}
For $a\in[K]$, let $(\muhat_{a,n})_{n\geq1}$ satisfy polynomial concentration.  Suppose the arms are sampled using
\[
a_n=\arg\max_a\left\{\muhat_{a,T_a(n)}+C_1\frac{n^{b/\beta}}{T_a(n)^{\alpha/\beta}}+\frac{C_2}{T_a(n)}\right\}.
\]
If either $1\leq p\leq2$ and $\alpha\leq\beta/2$, or $p>2$ and $0<\alpha-\beta/p<1$, and if $\alpha(1-b/\alpha)\leq b<\alpha$, then the power mean $\muhat_n(p)$ satisfies $\muhat_n(p)\xrightarrow{\alpha',\beta'}\mu^*$ with
\[
\alpha'=(b-1)(1-b/\alpha),\qquad \beta'=b-1.
\]
\end{theorem}

\begin{theorem}[MCTS-ER convergence]
\label{thm:main_mcts_er_statement}
Apply Algorithm~\ref{alg:mcts-er-main} with constants $\{b_i,\alpha_i,\beta_i\}_{i=0}^H$ satisfying the Stochastic-Power-UCT recursion and $b_i<\alpha_i$, $b_i>2$.  Then, for every node $s_h$,
\[
\Vhat_n(s_h)\xrightarrow{\alpha_h,\beta_h}\widetilde{V}(s_h),
\]
and for every internal node $s_h$ and action $a$,
\[
\Qhat_n(s_h,a)\xrightarrow{\alpha_{h+1},\beta_{h+1}}\widetilde{Q}(s_h,a).
\]
At the root, with the best parameter tuning,
\[
\left|\E[\Vhat_n(s_0)]-\widetilde{V}(s_0)\right|\leq O(n^{-1/2}).
\]
\end{theorem}
The proof is an induction over tree depth.  Leaf estimates concentrate by standard rollout concentration; Q-estimates inherit child-value concentration independently of the selection rule; and the root value follows from Theorem~\ref{thm:main_bandit_er_statement}.  The full argument is given in Appendix~\ref{app:full}.

\subsection{Convex Regularized MCTS-ER}
\label{sec:main_convex_results}
The convex-regularized proof follows the E3W analysis of MENTS/RENTS/TENTS.  The ER term is controlled through the uniform-mixing component: for large $t$, each action receives at least $\Omega(t/\log t)$ visits with high probability.  Consequently,
\begin{equation}
\frac{c_2}{T_{s,a}(t)}\leq \frac{4c_2\log(2+t)}{t},
\label{eq:main_er_perturb}
\end{equation}
with the same high-probability envelope as the visit-count concentration.  Since $\sum_{k\leq t}(\log k)/k=O((\log t)^2)=o(t/\log t)$, the cumulative ER perturbation is lower order in the E3W induction.

\begin{theorem}[Exponential value convergence of E3W-ER]
\label{thm:main_exp_value_statement}
At the root node $s$ with $N(s)$ visits, for any strongly convex regularizer $\Omega$ and any $\epsilon>0$,
\[
\Prob\left(|V_\Omega^{\mathrm{ER}}(s)-V_\Omega^*(s)|>\epsilon\right)
\leq C_{\mathrm{ER}}\exp\left\{-\frac{N(s)}{\hat C_{\mathrm{ER}}\sigma\log^2(2+N(s))}\right\}.
\]
The returned action also satisfies, for large $t$,
\[
\Prob(a_t\neq a^*)\leq C_{\mathrm{ER}}t\exp\left\{-\frac{t}{\hat C_{\mathrm{ER}}\sigma(\log t)^3}\right\}.
\]
\end{theorem}

\begin{theorem}[Regret and value error]
\label{thm:main_regret_error_statement}
Let $D_{\Omega^*}$ be the Bregman divergence of $\Omega^*$.  For E3W-ER, the expected pseudo-regret obeys
\[
\E[R_n]\leq -\tau\Omega(\hat\pi)+\sum_{t=1}^nD_{\Omega^*}(\hat V^t+V,\hat V^t)+O(n/\log n)+c_2H_n,
\]
where $H_n=\sum_{t=1}^n\sum_a\pi_t^{\mathrm{ER}}(a|s)/T_{s,a}(t)=O(|\mathcal A|\log n)$.  Moreover, with probability at least $1-\delta$,
\[
-\Psi_{\mathrm{ER}}-\frac{\tau(U_\Omega-L_\Omega^{\min})}{1-\gamma}
\leq V_\Omega^{\mathrm{ER}}(s)-V^*(s)
\leq \Psi_{\mathrm{ER}}+\frac{c_2}{\tau\min_aT_{s,a}},
\]
where $\Psi_{\mathrm{ER}}=\sqrt{\hat C_{\mathrm{ER}}\sigma^2\log(C_{\mathrm{ER}}/\delta)/(2N(s))}$.
\end{theorem}

\begin{table}[t]
\centering
\caption{Main ER-augmented bounds.  ER changes constants or lower-order terms, not leading asymptotic rates.}
\label{tab:main_summary}
\small
\begin{tabular}{lcccc}
\toprule
Algorithm & Conv. rate & Regret leading term & Error lower term & Sparse? \\
\midrule
Power-UCT-ER & $O(n^{-1/2})$ & -- & $O(n^{-1/2})$ & No \\
MENTS-ER & $\exp(-t/(\log t)^3)$ & $n|\mathcal A|/\tau$ & $\tau\log|\mathcal A|/(1-\gamma)$ & No \\
RENTS-ER & $\exp(-t/(\log t)^3)$ & $n|\mathcal A|/\tau$ & $\tau(\log|\mathcal A|-\log m)/(1-\gamma)$ & No \\
TENTS-ER & $\exp(-t/(\log t)^3)$ & $n|\mathcal K_{\mathrm{ER}}|/2$ & $\tau(|\mathcal A|-1)/(2|\mathcal A|(1-\gamma))$ & Yes \\
\bottomrule
\end{tabular}
\end{table}



\section{Proof Roadmap and Consequences}
\label{sec:main_proof_roadmap}
This section records the intermediate statements that connect the algorithms to the main rates.  The proof bodies are not rewritten here; they are kept in Appendix~\ref{app:full}.  The purpose of the restructured main text is to make the submission self-contained at the level of assumptions, rates, and dependencies.

\subsection{Power-mean ingredients}
The power-mean analysis begins with the combined confidence index
\[
U^{\mathrm{ER}}_{a,n,s}=\muhat_{a,s}+C_1\frac{n^{b/\beta}}{s^{\alpha/\beta}}+\frac{C_2}{s}.
\]
The following two facts are the only places where ER changes the original Stochastic-Power-UCT proof.

\begin{lemma}[Dominance of the ER bonus]
\label{lem:main_dominance_statement}
If $\alpha/\beta\leq1/2$, then $s^{-1}=o(s^{-\alpha/\beta})$.  Consequently, for fixed $n$, the ER term $C_2/s$ decays faster in the child visit count than the polynomial term $C_1n^{b/\beta}s^{-\alpha/\beta}$, and the combined bonus in \eqref{eq:main_combined_bonus} has the same leading order as the original polynomial bonus.
\end{lemma}

\begin{proposition}[ER-modified suboptimal-arm threshold]
\label{prop:main_threshold_statement}
For each suboptimal arm $a$ with gap $\Delta_a$, define
\[
A_a^{\mathrm{ER}}(n)=\inf\left\{s\leq n:C_1\frac{n^{b/\beta}}{s^{\alpha/\beta}}+\frac{C_2}{s}\leq\frac{\Delta_a}{2}\right\}.
\]
Then
\[
A_a^{\mathrm{ER}}(n)\leq \left(\frac{2C_1}{\Delta_a}\right)^{\beta/\alpha}n^{b/\alpha}+\frac{4C_2}{\Delta_a},
\qquad
A_a^{\mathrm{ER}}(n)=A_a(n)(1+o(1)).
\]
\end{proposition}

Given Proposition~\ref{prop:main_threshold_statement}, the optimal-arm concentration, suboptimal-arm count, and power-mean error decomposition are exactly the same chain as in Stochastic-Power-UCT with $A_a(n)$ replaced by $A_a^{\mathrm{ER}}(n)$.  In particular, for any $u\geq A_a^{\mathrm{ER}}(n)$,
\[
\Prob(T_a(n)\geq u)\leq \frac{2cC_1^{-\beta}}{b-1}(u-1)^{-(b-1)},
\]
and the power-mean error admits the decomposition
\[
|\muhat_n(p)-\mu^*|\leq R\sum_{a\neq a^*}\frac{T_a(n)}{n}
+\left(\sum_{a=1}^K\frac{T_a(n)}{n}|\muhat_{a,T_a(n)}-\mu_a|^p\right)^{1/p}.
\]
Thus ER only modifies constants through $C_2$ and the minimum gap $\Delta=\min_{a:\mu_a<\mu^*}\Delta_a$.

\subsection{Convex-regularized ingredients}
For E3W-ER, the regularized policy map is Lipschitz because $\Omega$ is strongly convex.  Hence
\[
\|\pi^{\mathrm{ER}}(\cdot|s)-\pi^*(\cdot|s)\|_p
\leq L\left\|\frac{Q^{\mathrm{ER}}(s,\cdot)-Q^*(s,\cdot)}{\tau}\right\|_p.
\]
The Q-perturbation decomposes into the original estimation error plus the ER vector:
\[
\left\|\frac{Q^{\mathrm{ER}}(s,\cdot)-Q^*(s,\cdot)}{\tau}\right\|_p
\leq \left\|\frac{Q(s,\cdot)-Q^*(s,\cdot)}{\tau}\right\|_p
+\frac{c_2}{\tau}\left\|\frac{1}{T_{s,\cdot}(t)}\right\|_p.
\]
The E3W uniform component gives $T_{s,a}(t)\geq t/(4\log(2+t))$ with high probability, which yields \eqref{eq:main_er_perturb}.  Summing this term across the induction contributes only $O((\log t)^2)$, while the dominant visit-count deviation is $O(t/\log t)$.  This separation is the reason the ER augmentation preserves the original exponential convergence envelope.

\begin{lemma}[Value perturbation]
\label{lem:main_value_perturb_statement}
For any strongly convex regularizer $\Omega$,
\[
|V_\Omega^{\mathrm{ER}}(s)-V_\Omega(s)|
=|\Omega^*(Q_\Omega^{\mathrm{ER}}/\tau)-\Omega^*(Q_\Omega/\tau)|
\leq \frac{c_2}{\tau\min_aT_{s,a}(t)}.
\]
In particular this perturbation is $O(\log t/t)$ under the E3W-ER visit lower bound.
\end{lemma}

\subsection{Regularizer-Specific Consequences}
\label{sec:main_regularizer_consequences}
The general regret and error bounds in Theorem~\ref{thm:main_regret_error_statement} become concrete after substituting entropy-specific Bregman and regularization constants.

\paragraph{MENTS-ER.}
For maximum entropy, $\Omega(\pi)=\sum_a\pi(a)\log\pi(a)$ and $\pi^{\mathrm{ER}}$ is a softmax over $Q+c_2/T$.  The regret and value-error bounds specialize to
\[
\E[R_n]\leq \tau\log|\mathcal A|+\frac{n|\mathcal A|}{\tau}+O(n/\log n)+O(c_2|\mathcal A|\log n),
\]
\[
-\Psi_{\mathrm{ER}}-\frac{\tau\log|\mathcal A|}{1-\gamma}
\leq V_{\mathrm{MENTS}}^{\mathrm{ER}}(s)-V^*(s)
\leq \Psi_{\mathrm{ER}}+O(\log t/t).
\]

\paragraph{RENTS-ER.}
For relative entropy, the ER term shifts the Q-values inside the KL-regularized update while retaining the previous-policy reference distribution.  With $m=\min_a\pi_{t-1}(a|s)$,
\[
\E[R_n]\leq \tau(\log|\mathcal A|-1/m)+\frac{n|\mathcal A|}{\tau}+O(n/\log n)+O(c_2|\mathcal A|\log n),
\]
\[
-\Psi_{\mathrm{ER}}-\frac{\tau(\log|\mathcal A|-\log(1/m))}{1-\gamma}
\leq V_{\mathrm{RENTS}}^{\mathrm{ER}}(s)-V^*(s)
\leq \Psi_{\mathrm{ER}}+O(\log t/t).
\]

\paragraph{TENTS-ER.}
For Tsallis entropy, $\Omega(\pi)=\frac12(\|\pi\|_2^2-1)$ and the active set is
\[
\mathcal K_{\mathrm{ER}}(s)=\left\{a:\frac{Q(s,a)+c_2/T_{s,a}}{\tau}>\psi\left(\frac{Q(s,\cdot)+c_2/T_{s,\cdot}}{\tau}\right)\right\}.
\]
The corresponding bounds are
\[
\E[R_n]\leq \frac{\tau(|\mathcal A|-1)}{|\mathcal A|}+\frac{n|\mathcal K_{\mathrm{ER}}|}{2}+O(n/\log n)+O(c_2|\mathcal A|\log n),
\]
\[
-\Psi_{\mathrm{ER}}-\frac{\tau(|\mathcal A|-1)}{2|\mathcal A|(1-\gamma)}
\leq V_{\mathrm{TENTS}}^{\mathrm{ER}}(s)-V^*(s)
\leq \Psi_{\mathrm{ER}}+O(\log t/t).
\]
The active set can be larger early in search because small $T_{s,a}$ raises the ER-augmented score of under-sampled actions.  As $T_{s,a}\to\infty$, $\mathcal K_{\mathrm{ER}}(s)$ approaches the original Tsallis active set, recovering sparse exploitation without manually annealing the regularizer.



\subsection{Technical Statements Kept in the Main Submission}
\label{sec:main_technical_statements}
The following statements are included in the main submission because they are the interfaces used by the final theorems.  Their proofs are in the appendix and are not edited.

\begin{table}[t]
\centering
\caption{Power-UCT-ER parameter conditions inherited from Stochastic-Power-UCT.}
\label{tab:main_conditions}
\small
\begin{tabular}{p{0.92\linewidth}}
\toprule
For every depth index $i\in\{0,\ldots,H\}$, choose constants satisfying all conditions below.\\
\midrule
$b_i<\alpha_i$ and $b_i>2$.\\
If $1\leq p\leq2$, require $\alpha_i\leq\beta_i/2$; if $p>2$, require $\alpha_i\leq\beta_i/2$ and $0<\alpha_i-\beta_i/p<1$.\\
$\alpha_i(1-b_i/\alpha_i)\leq b_i<\alpha_i$.\\
$\alpha_i=(b_{i+1}-1)(1-b_{i+1}/\alpha_{i+1})$ and $\beta_i=b_{i+1}-1$.\\
\bottomrule
\end{tabular}
\end{table}

\paragraph{Power-mean concentration chain.}
For a suboptimal arm $a$, the ER-modified UCB concentration has two sides:
\[
\Prob(U^{\mathrm{ER}}_{a,n,s}<\mu_a)\leq cC_1^{-\beta}n^{-b},
\qquad
\Prob(U^{\mathrm{ER}}_{a,n,s}>\mu^*)\leq cC_1^{-\beta}n^{-b}
\quad(s\geq A_a^{\mathrm{ER}}(n)).
\]
These inequalities imply a polynomial tail on the number of suboptimal pulls and then feed into the power-mean error decomposition.  The key point for the submission is that $A_a^{\mathrm{ER}}(n)=A_a(n)+O(C_2/\Delta_a)$, so the tree induction sees the same exponents $(\alpha_i,\beta_i)$ as before.  The final constants depend on the ER strength $C_2$, the number of actions, and the smallest positive gap, but not on a new asymptotic exponent.

\paragraph{Convex visit-count chain.}
For E3W-ER, let $N_t(s,a)$ be the actual visit count and $N_t^*(s,a)=\pi^*(a|s)t$ be the ideal count under the limiting regularized policy.  The appendix proves that, for constants depending on the regularizer, action set, temperature, and ER coefficient,
\[
\Prob\left(|N_t(s,a)-N_t^*(s,a)|>\frac{C_{\mathrm{ER}}t}{\log t}\right)
\leq \hat C_{\mathrm{ER}}t\exp\left(-\frac{t}{(\log t)^3}\right).
\]
Together with Lemma~\ref{lem:main_value_perturb_statement}, this yields the value convergence in Theorem~\ref{thm:main_exp_value_statement}.  The ER contribution is lower order because the cumulative perturbation along the induction is bounded by
\[
\sum_{k=1}^t\frac{\log(2+k)}{k}=O((\log t)^2),
\]
whereas the visit-count tolerance is $O(t/\log t)$.

\paragraph{What remains in the appendix.}
Appendix~\ref{app:full} contains the original proof of: (i) ER dominance and the ER-modified threshold; (ii) the modified analogues of the Stochastic-Power-UCT Lemmas 4--10; (iii) the induction proving convergence of every $\Vhat_n(s_h)$ and $\Qhat_n(s_h,a)$; (iv) leaf concentration, Lipschitz policy perturbation, expected and actual visit-count concentration for E3W-ER; and (v) the regret/error specialization for MENTS-ER, RENTS-ER, and TENTS-ER.  The proof structure is summarized here only to make the nine-page main paper readable.

\section{Experimental Results}
\label{sec:experiments}
\begin{figure}[h!]
    \centering
    \begin{minipage}{0.5\textwidth}
        \centering
        \includegraphics[width=\textwidth]{exp/fl812.png}
        \caption{Frozen Lake 8x12}
        \label{fig:fl812}
    \end{minipage}\hfill
    \begin{minipage}{0.5\textwidth}
        \centering
        \includegraphics[width=\textwidth]{exp/sailing66.png}
        \caption{Sailing 6x6}
        \label{fig:sailing66}
    \end{minipage}
\end{figure}

\section{Discussion}
\label{sec:main_discussion}
\subsection{Finite-Time Interpretation}
\label{sec:main_finite_time}
The ER term does not improve the leading asymptotic rate, but it changes how simulations are allocated before the asymptotic regime dominates.  The intended effect is most visible in sparse planning, where the simulation budget is small relative to the number of reachable tree edges.

\paragraph{Path-level uncertainty.}
A local UCB-style count at a deep node can be misleading when the node is reached through an unreliable prefix.  ER assigns the trajectory-level score $\sum_{e\in\mathrm{path}}1/N(e)$, so one weak prefix edge remains visible to all descendants.  This is the structural information absent from purely node-local bonuses.

\paragraph{Bottleneck detection.}
If a high-value branch depends on a rarely traversed edge, then the term $1/T_{s,a}$ keeps that edge attractive until it has enough support.  Once the edge is reinforced, the ER contribution decays faster than the polynomial bonus in Power-UCT-ER and at rate $O(\log t/t)$ under E3W-ER.  This explains why ER can change finite-time exploration without changing the limiting rate.

\paragraph{Complementarity with regularization.}
Entropy regularization smooths action probabilities according to Q-values, while ER modifies the Q-values according to structural uncertainty.  In MENTS-ER and RENTS-ER this produces a soft preference for under-sampled edges.  In TENTS-ER it can temporarily enlarge the active set $\mathcal K_{\mathrm{ER}}$, after which the policy returns to sparse exploitation as visits accumulate.

\paragraph{Coefficient choice.}
The constants $C_2$ and $c_2$ control the finite-time strength of the structural exploration signal.  Large coefficients may over-emphasize rarely visited edges, whereas small coefficients recover the base algorithms.  The theoretical results are therefore rate-preservation statements: ER is safe asymptotically, and the empirical value of the coefficient should be assessed in the reserved experimental page.

\subsection{Limitations}
\label{sec:main_discussion_structure}
ER is most useful in sparse planning regimes, where a small number of simulations must cover a large search tree.  The bonus detects bottleneck edges because $\Reff$ accumulates along the trajectory; it is therefore complementary to local value uncertainty in UCB-style selection and to entropy-based smoothing in regularized planning.  In a pure tree, this path-level signal reduces exactly to a local $1/T_{s,a}$ term at decision time, enabling the same node-wise bandit decomposition used by the base algorithms.

The extension to state graphs with transpositions is more delicate.  In a graph, adding a single edge can change effective resistances globally, so the local decomposition no longer holds.  A possible future proof route is to control a global potential such as $\Phi_n=\sum_s\Reff(s_0,s;G_n)$, which decreases under Rayleigh monotonicity as the discovered graph becomes more connected.  We leave this graph-level ER analysis for future work.

The main technical results are stated in Sections~\ref{sec:main_power_results}--\ref{sec:main_convex_results}. The original proof bodies and technical lemmas are placed after the references in Appendix~\ref{app:full}, without rewriting the proofs.

\paragraph{Implementation note.}
In a tree implementation, no matrix inverse or Laplacian solve is needed to use ER.  The effective-resistance path sum in \eqref{eq:main_er_path} can be maintained incrementally along each simulation, and the action-dependent term at selection time is exactly $1/T_{s,a}$.  Thus Power-UCT-ER has the same per-action selection complexity as the base polynomial-bonus rule, and E3W-ER only adds a scalar offset before the regularized policy is evaluated.  This matters for experiments because any finite-time gain should not be attributed to extra planning depth or an expensive global graph computation.

\paragraph{Limitations of the current draft.}
The theory is asymptotic and does not by itself identify the best value of $C_2$ or $c_2$ for a finite simulation budget.  Large ER coefficients can over-emphasize rarely visited edges, while very small coefficients recover the base algorithms.  The reserved experimental page should therefore include sensitivity to the ER coefficient or use a fixed, clearly stated tuning rule.  The present draft makes no empirical claim until those measurements are added.


% Reserved for future experimental results.
\clearpage

\bibliographystyle{plainnat}
\bibliography{references}

\clearpage
\appendix

\section{Full Technical Details and Proofs}
\label{app:full}
\input{appendix}

\clearpage
\input{checklist}
\end{document}
